% Standalone problem document, generated from the adjacent README.md.
% Compile with XeLaTeX. Mathematical content is maintained in Markdown.
\documentclass[11pt,a4paper]{article}
\usepackage[fontsize=10.5pt]{scrextend}
\usepackage[margin=22mm,headheight=15pt,headsep=9mm,footskip=11mm]{geometry}
\usepackage{fontspec}
\setmainfont{texgyrepagella}[Extension=.otf,UprightFont=*-regular,BoldFont=*-bold,ItalicFont=*-italic,BoldItalicFont=*-bolditalic]
\setsansfont{texgyreheros}[Extension=.otf,UprightFont=*-regular,BoldFont=*-bold,ItalicFont=*-italic,BoldItalicFont=*-bolditalic]
\setmonofont{texgyrecursor}[Extension=.otf,UprightFont=*-regular,BoldFont=*-bold,ItalicFont=*-italic,BoldItalicFont=*-bolditalic,Scale=MatchLowercase]
\usepackage{amsmath,amssymb}
\usepackage{unicode-math}
\setmathfont{texgyrepagella-math.otf}
\usepackage{xcolor,microtype,enumitem,needspace,fancyhdr}
\usepackage{bookmark}
\definecolor{ink}{HTML}{17344B}
\definecolor{accent}{HTML}{197D80}
\definecolor{muted}{HTML}{596773}
\hypersetup{colorlinks=true,linkcolor=accent,urlcolor=accent,pdftitle={RA-14 - Optimal
query complexity of spectral rank-k
approximation},pdfauthor={Open Problems in Numerical Linear Algebra}}
\urlstyle{same}
\setlength{\parindent}{0pt}
\setlength{\parskip}{5pt plus 1pt minus 1pt}
\linespread{1.04}
\setlength{\emergencystretch}{3em}
\widowpenalty=10000
\clubpenalty=10000
\displaywidowpenalty=10000
\allowdisplaybreaks[1]
\setlist{nosep,leftmargin=1.5em}
\providecommand{\tightlist}{\setlength{\itemsep}{0pt}\setlength{\parskip}{0pt}}
\providecommand{\pandocbounded}[1]{#1}
\pagestyle{fancy}
\fancyhf{}
\fancyhead[L]{\sffamily\footnotesize\color{muted}OPEN PROBLEMS IN NUMERICAL LINEAR ALGEBRA}
\fancyhead[R]{\sffamily\bfseries\footnotesize\color{accent}RA-14}
\fancyfoot[L]{\parbox[b]{0.9\textwidth}{\sffamily\fontsize{8}{10}\selectfont\color{muted}Editorial ratings; a bounded search cannot certify that a problem remains unsolved.\\Literature check: 2026-09-13}}
\fancyfoot[R]{\sffamily\footnotesize\color{muted}\thepage}
\renewcommand{\headrulewidth}{0.3pt}
\renewcommand{\headrule}{\hbox to\headwidth{\color{accent}\leaders\hrule height \headrulewidth\hfill}}
\makeatletter
\renewcommand{\section}{\Needspace{5\baselineskip}\@startsection{section}{1}{0pt}{12pt plus 2pt minus 2pt}{4pt}{\sffamily\bfseries\large\color{ink}}}
\renewcommand{\subsection}{\Needspace{4\baselineskip}\@startsection{subsection}{2}{0pt}{9pt plus 1pt minus 1pt}{3pt}{\sffamily\bfseries\normalsize\color{ink}}}
\makeatother
\setcounter{secnumdepth}{0}
\begin{document}
{\sffamily\small\color{muted}Randomized and low-rank approximation\par}
\vspace{3pt}
{\raggedright\hyphenpenalty=10000\exhyphenpenalty=10000\sffamily\bfseries\fontsize{21}{24}\selectfont\color{ink}Optimal
query complexity of spectral rank-\(k\) approximation\par}
\vspace{7pt}
{\sffamily\small\textbf{Difficulty:} extreme\quad\textbf{Importance:} broadly
interesting\par}
{\sffamily\small\bfseries\color{accent}Status: Partially resolved\par}
\vspace{4pt}
{\color{accent}\hrule height 0.8pt}
\vspace{6pt}
\textbf{Rating rationale:} Extreme because matching information bounds
for every adaptive algorithm and growing rank is a fundamental barrier;
broad impact includes large-scale spectral computation and data
analysis.\\
\textbf{Source:} Bakshi--Narayanan, Open Question 1.10.

\subsection{Problem statement}\label{problem-statement}

Let \(n\ge2\), \(1\le k< n\), and \(0<\varepsilon<1/2\). An unknown
matrix \(A\in\mathbb R^{n\times n}\) is available through exact products
\(Ax\) and \(A^Tx\), with one such product counting as one query. An
algorithm may choose each real query vector adaptively using all earlier
answers and its random bits. Arithmetic and other computation between
queries are unrestricted; matrix entries are not otherwise accessible.

Let \(q_{\mathrm{sp}}(n,k,\varepsilon)\) be the smallest integer \(q\)
for which an algorithm using at most \(q\) queries returns a matrix
\(Z\in\mathbb R^{n\times k}\) with \(Z^TZ=I_k\) such that, for every
input \(A\),

\[\Pr\!\left[
\|A(I-ZZ^T)\|_2
\le (1+\varepsilon)
\min_{\substack{U\in\mathbb R^{n\times k}\\U^TU=I_k}}
\|A(I-UU^T)\|_2
\right]\ge\frac{99}{100}.\]

Probability is over the algorithm's internal randomness, and
\(\|\cdot\|_2\) is the spectral norm.

\textbf{Open problem.} Determine \(q_{\mathrm{sp}}(n,k,\varepsilon)\) up
to universal constant factors, with the dependence on \(k\), \(n\), and
\(\varepsilon\) simultaneous. In particular, determine the optimal
dependence on a growing target rank rather than keeping \(k\) fixed. The
answer must cover the finite-dimensional regime, where learning all
columns uses \(n\) queries.\\
This is a square-matrix formulation of the source's question with an
editorially fixed two-sided oracle model. The introduction describes
products \(Av\); Algorithm 7.4 explicitly uses both \(A\) and \(A^T\).
On symmetric inputs these models coincide.

\subsection{Why it matters in numerical linear
algebra}\label{why-it-matters-in-numerical-linear-algebra}

Matrix products dominate many large-scale singular-subspace
computations. The question asks which dependence on the requested number
of singular vectors is unavoidable for every adaptive algorithm, rather
than for a chosen Krylov implementation.

\newpage
\subsection{References}\label{references}

\begin{enumerate}
\def\labelenumi{\arabic{enumi}.}
\tightlist
\item
  Ainesh Bakshi and Shyam Narayanan,
  \href{https://arxiv.org/html/2304.03191v1\#S1.SS2}{\emph{Krylov
  Methods are (nearly) Optimal for Low-Rank Approximation}},
  arXiv:2304.03191v1 (2023), Open Question 1.10; Theorem 1.1 and
  Algorithm 7.4.
\item
  Tyler Chen et al., \href{https://arxiv.org/abs/2508.06486}{\emph{Does
  block size matter in randomized block Krylov low-rank
  approximation?}}, 2025, §1: later block-size bounds.
\end{enumerate}

\subsection{Status check}\label{status-check}

The source's \(O(k\log n/\sqrt\varepsilon)\) upper bound and fixed-rank
lower bound do not determine the growing-rank dependence. The 2025 paper
studies Krylov block sizes; its clustered-gap conjecture is separately
cataloged.

On 2026-09-08 the \href{https://arxiv.org/abs/2304.03191}{source record}
still listed only v1. Searches for ``Open Question 1.10,'' target-rank
matrix-vector complexity, and 2025/2026 follow-ups found no joint
characterization. The source's lower bound has a
sufficiently-large-dimension regime, not every finite \(n,\varepsilon\).
This is a bounded check.

\subsection{Audit --- 2026-09-10}\label{audit-2026-09-10}

Rechecked \href{https://arxiv.org/html/2304.03191v1}{Bakshi--Narayanan,
Theorem 1.1 and Open Question 1.10}. Fixed-rank spectral complexity is
settled in the theorem's sufficiently-large-dimension regime;
growing-rank and simultaneous finite-parameter dependence remain
unresolved. Later query-complexity searches and the
\href{https://doi.org/10.1137/1.9781611978971.42}{SODA 2026 block-size
paper} did not settle the full target.

\subsection{Finite-accuracy partial result --- 13 September
2026}\label{finite-accuracy-partial-result-13-september-2026}

\textbf{Author:} Sidney Holden, Center for Computational Biology,
Flatiron Institute, Simons Foundation.
\href{https://github.com/ajt60gaibb/OpenProblemsInNLA/blob/main/references/holden-ra14-v5-2026-09-13/README.md}{Submission
and verified affiliation}.

\href{https://github.com/ajt60gaibb/OpenProblemsInNLA/blob/main/references/holden-ra14-v5-2026-09-13/report.pdf}{Theorem
1.1} establishes, for a universal positive constant \(c\) and every
original admissible parameter triple,

\[q_{\mathrm{sp}}(n,k,\varepsilon)\ge c\frac{k}{\sqrt{\varepsilon}}\log\left(1+\frac{n\sqrt{\varepsilon}}{k}\right).\]

Together with the reproduced upper bound, this gives matching
universal-factor bounds when \(\varepsilon\le(k/n)^2\) and when
\(\varepsilon\ge k/n\), including
\(q_{\mathrm{sp}}(n,1,1/n)=\Theta(\sqrt n\log n)\) for \(n>2\). Sections
3--10 prove the new lower bound; Appendices A and B reproduce the
charged reduction and upper bound.
\href{https://github.com/ajt60gaibb/OpenProblemsInNLA/blob/main/references/holden-ra14-v5-2026-09-13/report.tex}{Proof
source}.

The stated partial result passed a separate
\href{https://github.com/ajt60gaibb/OpenProblemsInNLA/blob/main/references/holden-ra14-v5-2026-09-13/independent-review.md}{independent
informal Codex AI-agent audit}. \textbf{RA-14 remains Partially
resolved:} at \(k=1\) and \(\varepsilon=(\log n/n)^2\), the bounds still
leave \(\Omega(n\log\log n/\log n)\) versus \(O(n)\), an unbounded
factor. The full simultaneous universal-factor target is unchanged. No
Lean verification, external human peer review or priority claim is
asserted. This continues the earlier partial submission
\href{https://github.com/ajt60gaibb/OpenProblemsInNLA/pull/190}{PR
\#190}; it does not duplicate a previously pushed full solution.
\end{document}
